% ============================================================================
%  SECCIÓN 4.2: CREACIÓN DE MODELOS DE DATOS
% ============================================================================

\section{Creaci\'on de modelos de datos}

Los modelos de datos definen la estructura de la informaci\'on que maneja la
aplicaci\'on. En un proyecto con React Native y TypeScript, los modelos se
implementan como interfaces o tipos que permiten el tipado est\'atico y la
detecci\'on temprana de errores durante el desarrollo.

\subsection{Interfaces principales}

Los modelos de datos del proyecto se definen en el directorio
\textit{src/types/} y se organizan de acuerdo con las entidades del dominio
del sistema.

\subsubsection{Modelo de escenario}

El escenario deportivo es la entidad central del sistema. Cada escenario
contiene informaci\'on b\'asica, una relaci\'on con deportes asociados, im\'agenes y
eventos:

\begin{lstlisting}
export interface Scenario {
  id: string;
  nombre: string;
  tipo: string;
  direccion: string;
  capacidad: number;
  descripcion: string | null;
  estado: 'activo' | 'inactivo';
  latitud: number | null;
  longitud: number | null;
  created_at: string;
}
\end{lstlisting}

Para las pantallas de detalle y cat\'alogo, se extiende el modelo con las
relaciones necesarias:

\begin{lstlisting}
export interface ScenarioWithDetails extends Scenario {
  scenario_images: ScenarioImage[];
  scenario_sports: ScenarioSport[];
  events: Event[];
}
\end{lstlisting}

\subsubsection{Modelo de imagen del escenario}

Cada escenario puede tener m\'ultiples im\'agenes, una de las cuales se marca
como primaria:

\begin{lstlisting}
export interface ScenarioImage {
  url: string;
  is_primary: boolean;
  storage_path: string;
  display_order: number;
}
\end{lstlisting}

\subsubsection{Modelo de deporte asociado}

La relaci\'on entre escenarios y deportes se modela como una tabla intermedia
que permite una relaci\'on de muchos a muchos:

\begin{lstlisting}
export interface ScenarioSport {
  sports: { nombre: string };
}
\end{lstlisting}

\subsubsection{Modelo de evento}

Los eventos representan actividades pr\'oximas en cada escenario:

\begin{lstlisting}
export interface Event {
  id: string;
  scenario_id: string;
  nombre: string;
  descripcion: string | null;
  fecha: string;
  hora: string;
  created_at: string;
}
\end{lstlisting}

\subsubsection{Modelo de favorito}

El modelo de favorito establece la relaci\'on entre un usuario y un escenario
guardado:

\begin{lstlisting}
export interface Favorite {
  user_id: string;
  scenario_id: string;
  created_at: string;
}
\end{lstlisting}

\subsection{Resumen de modelos}

\begin{table}[H]
  \centering
  \caption{Modelos de datos del sistema}
  \contenidoConFuente{%
    \begin{tablaAPA}
      \begin{tabular}{@{}p{3.2cm}p{5.0cm}p{7.4cm}@{}}
        \toprule
        \textbf{Modelo} & \textbf{Campos principales} & \textbf{Prop\'osito} \\
        \midrule
        \textit{Scenario} & id, nombre, tipo, direccion, capacidad,
        estado & Entidad central: escenario deportivo. \\
        \addlinespace
        \textit{ScenarioImage} & url, is\_primary, display\_order &
        Im\'agenes asociadas a cada escenario. \\
        \addlinespace
        \textit{ScenarioSport} & sports (nombre) & Deportes disponibles
        en el escenario. \\
        \addlinespace
        \textit{Event} & id, nombre, fecha, hora, descripcion &
        Eventos pr\'oximos en cada escenario. \\
        \addlinespace
        \textit{Favorite} & user\_id, scenario\_id & Escenarios guardados
        por el usuario. \\
        \bottomrule
      \end{tabular}
    \end{tablaAPA}%
  }{Elaboraci\'on propia en base a los archivos del directorio
  \textit{src/types/}.}
\end{table}
